\section{Counterfactual Design and Results}

The canonical submission pipeline searches over a small ladder of common shock sizes and keeps
the largest value that converges in both timing scenarios. The selected common shock size is
$\Delta=\SelectedDelta$. This choice reflects the assignment's emphasis on a working dynamic
model rather than on forcing a larger but unstable shock through the solver.

\input{generated/table_ai_results.tex}

The table shows the intended timing contrast. In the immediate scenario, the focus labor market
experiences positive early real-wage effects right away. In the anticipated scenario, the early
window is close to zero by construction, while the middle and late windows pick up once the
shock activates. That difference is the core anticipation result in this simplified exercise.

\begin{figure}[htbp]
    \centering
    \includegraphics[width=0.48\textwidth]{figures/dynamics_immediate.pdf}
    \includegraphics[width=0.48\textwidth]{figures/dynamics_anticipated.pdf}
    \caption{Dynamic adjustment under immediate and anticipated AI productivity shocks.}
    \label{fig:dynamics}
\end{figure}

\begin{figure}[htbp]
    \centering
    \includegraphics[width=0.75\textwidth]{figures/state_map_immediate.pdf}
    \caption{Cross-state real-wage effects in the immediate AI scenario.}
    \label{fig:map}
\end{figure}

Overall, the exercise should be interpreted as a compact CDP-style application rather than a
full structural study of AI. Its value is that it delivers a coherent dynamic counterfactual
with explicit timing, validated baseline mechanics, and a reproducible submission package.
